\documentclass[11pt]{letter}
\usepackage[a4paper,margin=1in]{geometry}
\usepackage{hyperref}
\hypersetup{colorlinks=true, urlcolor=blue, linkcolor=blue}

\signature{Quang-Vinh Dang}
\address{British University Vietnam \\ Hung Yen, Vietnam \\ vinh.dq4@buv.edu.vn \\ ORCID: 0000-0002-3877-8024}

\begin{document}

\begin{letter}{Editor-in-Chief \\ INFORMS Journal on Computing}

\opening{Dear Editor,}

I am pleased to submit my manuscript, ``\textit{Adaptive Influence Minimization via Deferred Frontier Blocking: Theory, Algorithms, and Exact Validation on Bounded-Pathwidth Graphs},'' for consideration for publication in \textit{INFORMS Journal on Computing}.

The manuscript studies influence minimization (IMIN): given a seeded rumor or epidemic on a network and a fixed budget of vertices that can be removed, which vertices should be chosen to minimize the eventual spread. Every published IMIN algorithm to date is one-shot, committing the entire budget before the cascade is observed even once. This paper makes two contributions that address gaps this raises:

\begin{itemize}
\item \textbf{DEFER}, a general adaptive wrapper that upgrades any one-shot IMIN planner into a round-by-round adaptive one. I prove that the pure-deferral variant is never worse than the one-shot plan it wraps; that the natural push-down commitment rule is the local myopic optimum but, by a closed-form counterexample, does not extend this domination guarantee to the full algorithm; and that the adaptivity gap of IMIN is unbounded, in contrast to the bounded adaptivity gap known for the complementary influence maximization problem.
\item \textbf{TWIG}, an exact dynamic program that computes the true expected spread of a candidate blocking set rather than estimating it by Monte Carlo simulation, on graphs of bounded pathwidth. I use TWIG to give the first simulation-noise-free measurement of how often and how badly greedy IMIN selection departs from the true optimum.
\end{itemize}

Both algorithms are implemented, empirically validated on real networks, and released in full at \url{https://github.com/vinhqdang/trace_e}, together with all data and benchmark scripts used in the paper.

This manuscript is original, has not been published previously, and is not under consideration at any other journal. I have no conflicts of interest to disclose. I believe the paper's combination of new theoretical guarantees for adaptive network intervention with an exact validation tool for the simulation-based evaluations the literature relies on is a good fit for the journal's scope in algorithms, optimization, and computational methods for networked systems.

Thank you for considering my manuscript. I look forward to your response.

\closing{Sincerely,}

\end{letter}
\end{document}
